\documentclass[12pt]{article}
\usepackage{amsmath}
\usepackage{multicol}
\usepackage[a4paper, margin=1in]{geometry}

\begin{document}

\noindent \Large \textbf{Dapatkan himpunan penyelesaian dari}
\normalsize

\begin{multicols}{2}
\begin{enumerate}
    \item $7x - 2 \leq 9x + 3$
    \item $-4 < 3x + 2 < 5$
    \item $-3 < 1 - 6x \leq 4$
    \item $x^2 + 2x - 12 < 0$
    \item $2x^2 + 5x - 3 > 0$
    \item $\displaystyle \frac{1}{3x-2} \leq 4$
    \item $\displaystyle \frac{x+9}{x-3} \leq 0$
    \item $\displaystyle \frac{3}{x+5} > 2$
    \item $x^4 - 2x^2 \geq 8$
    \item $x^2 - 3x - 4 \geq 0$
    
    \columnbreak

    \setcounter{enumi}{10}
    \item $|2x+1|-x=|x+2|$
    \item $|x-1|+|x-2|=1$
    \item $|4x + 5| \leq 10$
    \item $|2x - 1| > 2$
    \item $\displaystyle \left| \frac{x}{4} + 1 \right| < 1$
    \item $|x - 1| < 2|x - 3|$
    \item $|2x - 1| \geq |x + 1|$
    \item $\displaystyle \left| \frac{x+2}{2x+1} \right| < 2$
    \item $|2x - 5| < 3x$
    \item $|3 - 7x| > 5 - x$
\end{enumerate}
\end{multicols}

\vspace{0.5cm}
\noindent \Large \textbf{Gambarkan sketsa grafik dari persamaan}
\normalsize
\begin{multicols}{2}
    \begin{enumerate}
        \item $y=(x+1)^2$
        \item $y=(x-2)^3-1$
        \item $y=|x+3|$
        \item $y=|x-1|+2$
        \item $y=\dfrac{1}{x+2}-3$
        \item $y=\sqrt{x+2}$
        \item $x=y^2-4$
        \item $x=(y-1)^2+1$
        \item $y=x^2-4x+3$
        \item $y=2x^2-3x+1$
    \end{enumerate}
    
\end{multicols}
\vspace{0.5cm}
\noindent \Large \textbf{Latihan Soal Persamaan Garis}
\normalsize
\begin{enumerate}
    \item Dapatkan $p$ dan $q$ jika diketahui garis yang melalui $(0,0)$ dan $(p,q)$ mempunyai gradien $\frac{1}{3}$, dan garis yang melalui $(5,4)$ dan $(p,q)$ mempunyai gradien $3$.
    
    \item Dapatkan koordinat semua titik $Q$ pada sumbu-$y$ sehingga garis yang melalui $M(-1,3)$ dan $Q$ tegak lurus ke garis yang melalui $N(5,-1)$ dan $Q$.
    
    \item Dapatkan gradien garis dan titik potong dengan sumbu-$y$ dari persamaan garis berikut ini:
    \begin{enumerate}
        \item $5x - 2y = 10$
        \item $\displaystyle \frac{x}{2} + \frac{y}{5} = 1$
    \end{enumerate}
    
    \item Gambarkan grafik untuk masing-masing persamaan yang diberikan berikut ini.
    \begin{enumerate}
        \item $\displaystyle \frac{x}{4} - \frac{y}{2} = 1$
        \item $\displaystyle x = \frac{1}{2}y + 3$
    \end{enumerate}
    
    \item Dapatkan nilai $m$, jika titik $(2, m)$ terletak pada garis yang melalui titik $(-1, 4)$ dan mempunyai gradien $-\frac{2}{3}$.
    
    \item Dapatkan koordinat $(u,v)$ jika diketahui garis yang melalui $(2,0)$ dan $(u,v)$ mempunyai gradien $2$, dan garis yang melalui $(-1,3)$ dan $(u,v)$ mempunyai gradien $-\frac{1}{2}$.
    
    \item Diketahui garis $g_1$ dengan gradien $m_1 = -\frac{2}{3}$, dan garis $g_2$ berpotongan tegak lurus dengan $g_1$ di titik $(-2, 3)$. Dapatkan persamaan garis $g_2$.
    
    \item Diketahui garis $k_1$ melalui titik $(-2, 4)$ dan $(2, -1)$. Jika $k_2$ adalah garis yang melalui titik $(3, 2)$ dan sejajar dengan garis $k_1$. Dapatkan persamaan garis $k_2$.
    
    \item Jika jarak antara titik $(p, -1)$ dengan titik $(2, 3)$ adalah $5$ maka tentukan semua nilai $p$.
    
    \item Jika titik tengah antara $P(-3, 2)$ dan titik $Q$ adalah titik $R(1, -2)$ maka tentukan koordinat titik $Q$.
    
    \item Diberikan titik koordinat $K(1, -2)$, $L(3, 1)$, dan $M(-2, 4)$. Jika garis $g$ melalui titik $K$ dan $L$ maka tentukan
    \begin{enumerate}
        \item jarak titik $M$ ke garis $g$.
        \item persamaan garis yang melalui $M$ dan sejajar dengan $g$.
    \end{enumerate}
    
    \item Tentukan gradien dari persamaan garis lurus yang melalui titik $(-2, 3)$ dan titik tengah antara titik $(4, -2)$ dan $(-2, 6)$.
    
    \item Tentukan persamaan garis $h$ yang melalui titik $(-1, 2)$ jika diketahui bahwa garis $h$ tegak lurus dengan $2x - cy = 3$ dan sejajar dengan $c^2x + 2(y+1) = 0$.
\end{enumerate}

\vspace{0.5cm}

\noindent \Large \textbf{Latihan Soal Persamaan Lingkaran dan Persamaan Parabola}
\normalsize

\begin{enumerate}
    \item Tentukan persamaan lingkaran yang memenuhi syarat berikut:
    \begin{enumerate}
        \item Pusat $(3, -2)$ dan berjari-jari $4$.
        \item Pusat $(-4, 5)$ dan menyinggung sumbu-$x$.
        \item Pusat $(2, 6)$ dan menyinggung sumbu-$y$.
    \end{enumerate}
    
    \item Tentukan persamaan lingkaran jika diketahui:
    \begin{enumerate}
        \item Garis tengahnya mempunyai ujung $(2, -1)$ dan $(6, 5)$.
        \item Pusat $(-1, 4)$ dan lingkaran melalui titik $(3, 1)$.
    \end{enumerate}

    \item Apakah persamaan di bawah ini merupakan lingkaran, titik, atau tidak memiliki grafik? Jika merupakan persamaan lingkaran, tentukan pusat dan jari-jarinya.
    \begin{enumerate}
        \item $x^2 + y^2 - 4x + 6y - 12 = 0$
        \item $x^2 + y^2 - 10x + 25 = 0$
        \item $3x^2 + 3y^2 - 6x + 12y = 0$
        \item $2x^2 + 2y^2 + 8x + 4y + 2 = 0$
    \end{enumerate}

    \item Dapatkan persamaan dari:
    \begin{enumerate}
        \item setengah bagian yang bawah dari lingkaran $x^2 + y^2 = 25$.
        \item setengah bagian yang atas dari lingkaran $x^2 + y^2 - 4x + 6y + 4 = 0$.
    \end{enumerate}

    \item Dapatkan persamaan dari:
    \begin{enumerate}
        \item setengah bagian yang kanan dari lingkaran $x^2 + y^2 = 36$.
        \item setengah bagian yang kiri dari lingkaran $x^2 + y^2 + 6x + 5 = 0$.
    \end{enumerate}

    \item Gambarkan grafik:
    \begin{enumerate}
        \item $y = \sqrt{16 - x^2}$
        \item $y = \sqrt{8 - 2x - x^2}$
    \end{enumerate}

    \item Gambarkan grafik:
    \begin{enumerate}
        \item $x = -\sqrt{9 - y^2}$
        \item $x = 2 + \sqrt{16 - y^2}$
    \end{enumerate}

    % Melompat ke nomor 13
    \item Suatu titik $(x,y)$ bergerak sedemikian hingga jaraknya ke $(0, 4)$ adalah $\sqrt{3}$ kali jaraknya ke $(0, 2)$.
    \begin{enumerate}
        \item Tunjukkan bahwa titik tersebut bergerak sepanjang suatu lingkaran.
        \item Dapatkan pusat dan jari-jarinya.
    \end{enumerate}

    \item Suatu titik $(x,y)$ bergerak sedemikian hingga jumlah kuadrat jaraknya ke $(3, 2)$ dan $(-1, -4)$ adalah $50$.

    % Melompat ke nomor 16
    \item Gambarkan grafik parabola dan tandai koordinat puncak serta perpotongannya dengan sumbu-sumbu koordinat:
    \begin{enumerate}
        \item $y = x^2 + 3$
        \item $y = x^2 - 4x + 3$
        \item $x = y^2 - 4y + 3$
    \end{enumerate}

    \item Dapatkan persamaan dari:
    \begin{enumerate}
        \item setengah bagian yang kanan dari parabola $y = 4 - x^2$.
        \item setengah bagian yang kiri dari parabola $y = x^2 + 4x$.
    \end{enumerate}

    \item Dapatkan persamaan dari:
    \begin{enumerate}
        \item setengah bagian yang atas dari parabola $x = y^2 - 4$.
        \item setengah bagian yang bawah dari parabola $x = y^2 - 2y - 3$.
    \end{enumerate}
\end{enumerate}


\end{document}